\section{Task × Timescale：科学超智能不是一个分数}

要理解为什么课程选择多智能体系统，先要区分两类能力：一种是在边界清楚的任务上极快、极准地完成计算；另一种是在多年尺度上穿越模糊目标、冲突证据和昂贵实验。前者可以产生“局部超人”，后者才接近讲者所说的 scientific superintelligence。

\subsection{先读坐标轴，再读模型名字}

图的纵轴从文献综述、实验、数据分析逐步上升到论文、博士课题、重大理论和范式突破，表示任务复杂度与影响；横轴从分钟延伸到一生，表示科学家完成任务所需时间。这个坐标不是精确测量，而是一种系统设计地图：越向右上，越需要长期记忆、资源调度、分解、验证和社会协作。

\lecturefigure{slide-004.jpg}{AlphaFold 在 Task × Timescale 坐标中的位置}{Stanford CS25 V6 Lecture 07 官方录像 00:13:10}

\begin{knowledgebox}{读图：AlphaFold 压缩的是一类专门任务}
AlphaFold 把蛋白质结构预测从可能需要多年实验推进的瓶颈，压缩为更短的计算流程，因此在横轴上产生巨大位移。但它解决的是定义明确、输入输出结构明确的专门问题。图把它放在右下方，是为了强调“超人级价值”可以来自一个窄任务的时间压缩，而不等于系统已经能提出完整研究议程。
\end{knowledgebox}

\teachervoice{00:10:17--00:12:45，老师用 AlphaFold 说明：一个系统可以在某个科学任务上产生革命性影响，却仍不具备开放式假设生成、跨领域综合和长期研究管理能力。}

\subsection{通用 LLM 覆盖了哪里，又缺了哪里}

课程随后把 Gemini、GPT、Claude 放在分钟到小时、文献综述到初步实验分析附近。它们的覆盖面比 AlphaFold 广，但大多停留在较短、较低风险的任务。这个图并不是说 LLM 只能做简单工作，而是在追问：怎样把通用知识与语言推理，延伸到需要几天、几月甚至多年验证的科学计划？

\lecturefigure{slide-005.jpg}{加入通用 LLM 后的 Task × Timescale 完整状态}{Stanford CS25 V6 Lecture 07 官方录像 00:15:54}

\begin{knowledgebox}{读图：先比较“宽度”，再比较“高度”}
AlphaFold 是窄而深：专门任务上的时间压缩非常大。通用 LLM 是宽而浅：能覆盖许多知识工作，却常缺少长期状态、稳定验证和实验闭环。AI co-scientist 的目标不是把某个气泡简单拖到右上角，而是组合模型、搜索、工具、记忆和人类实验，使系统能在更长时间、更高复杂度下保持可追踪进展。
\end{knowledgebox}

\begin{knowledgebox}{术语消化：三个常被混用的层级}
\begin{tabular}{L{0.16\textwidth}L{0.31\textwidth}L{0.36\textwidth}}
\toprule
层级 & 解决的问题 & 本讲中的例子 \\
\midrule
Task model & 在固定输入输出上优化一个专门任务 & AlphaFold 结构预测 \\
Foundation model & 在广泛任务中提供知识、语言与推理能力 & Gemini、GPT、Claude \\
Scientific system & 维持长期目标，组织搜索、批判、工具与验证 & AI co-scientist 多智能体系统 \\
\bottomrule
\end{tabular}
\end{knowledgebox}

\subsection{从 AlphaGo 学到的是搜索组织，而不是把科学变成棋盘}

AlphaGo 的启发在于：强策略网络之外，还可以用 self-play、搜索与比较持续改进候选。科学没有固定棋盘、完全规则和即时胜负信号，因此不能直接移植 Monte Carlo Tree Search；可移植的是\textbf{让多个候选互相施压、暴露弱点并迭代演化}的组织思想。

若一次系统运行生成 $N_0$ 个原始假设，经过 $K$ 轮比较与演化，每轮保留比例为 $\rho_k$、扩展比例为 $\gamma_k$，粗略的候选规模可写成
$$
N_{k+1}=\rho_kN_k+\gamma_kN_k.
$$
其中：
\begin{itemize}
  \item $N_k$ 是第 $k$ 轮仍被追踪的候选数量；
  \item $\rho_k$ 表示保留高价值候选的强度；
  \item $\gamma_k$ 表示从已有候选组合、变异或扩展出的新候选比例；
  \item 公式只描述搜索宽度，不代表科学质量会随 $N_k$ 单调上升。
\end{itemize}

\begin{warningbox}{科学不是有终局奖励的游戏}
棋类可以从胜负回传清晰信号；开放科学往往多年后才知道假设是否成立，而且“新颖”“可行”“重要”可能彼此冲突。把 tournament 引入科学，只能提供内部优先级，不能把 Elo 变成自然规律。
\end{warningbox}

\subsection{Scientific superintelligence 是协作结构}

前一节已经说明科学没有棋类那样明确的终局奖励，因此本节转向“谁来维持长期搜索”这一系统问题。课程用人脑与多个协作者的视觉，表达一种系统观：高水平科学不是单一模块瞬间给出答案，而是多条思路并行、彼此批判、借助工具和记忆，最终由人类把候选接入真实实验。Agent 的价值来自分工与反馈结构，而不只是复制更多相同模型实例。

\lecturefigure{slide-006.jpg}{科学超智能被表达为多智能体协作系统}{Stanford CS25 V6 Lecture 07 官方录像 00:18:36}

\begin{importantbox}{系统级能力的五个条件}
要从短任务走向长期科学，系统至少需要：持续目标与状态；并行探索与去重；显式反驳和不确定性；外部工具与证据；人类对实验、伦理和资源的最终控制。缺少其中任意一项，增加 agent 数量都可能只放大重复、错误或审查负担。
\end{importantbox}

\teachervoice{00:17:30--00:19:28，老师把目标描述为一组能够长期协作的 agent，而不是一个“知道所有事情”的模型。这个表述把架构问题放在模型能力之上。}

\subsection{本章小结}

Task × Timescale 图把两类进步分开：专门模型可以极大压缩一个科学任务，通用模型可以覆盖更多短任务，但开放式科学需要长期系统。AlphaGo 提供的是搜索、比较与演化的组织灵感，而不是科学评价函数。Scientific superintelligence 因而应被理解为模型、Agent、工具、记忆、人类专家和实验资源共同构成的协作结构。

